% Generated draft: requires mathematical and source review.
\documentclass[11pt]{article}
\usepackage[margin=25mm]{geometry}
\usepackage{fontspec}
\setmainfont{Latin Modern Roman}
\usepackage{amsmath,amssymb,mathtools,unicode-math}
\setmathfont{Latin Modern Math}
\usepackage{enumitem,needspace,xurl,hyperref}
\setlength{\parindent}{0pt}
\setlength{\parskip}{5pt}
\setlength{\emergencystretch}{3em}
\newcommand{\N}{\mathbb{N}}
\newcommand{\Z}{\mathbb{Z}}
\newcommand{\Q}{\mathbb{Q}}
\newcommand{\R}{\mathbb{R}}
\newcommand{\C}{\mathbb{C}}
\newcommand{\F}{\mathbb{F}}
\newcommand{\A}{\mathbb{A}}
\newcommand{\PP}{\mathbb{P}}
\newcommand{\E}{\mathbb{E}}
\newcommand{\eps}{\varepsilon}
\newcommand{\dd}{\,\mathrm{d}}
\begin{document}
\section*{UnsolvedMath ID: 30001508}
\textbf{Source problem code:} catalogue entry
\subsection{Definitions and mathematical statement}
A finite-sheeted cover of a manifold \(M\) is a covering map \(p\colon \widetilde M\to M\) with finitely many sheets. A manifold \(\widetilde M\) fibers over the circle if there is a locally trivial fiber bundle \(f\colon\widetilde M\to S^1\), where \(S^1=\{z\in\C:|z|=1\}\), whose fiber is a surface.

Under the standard orientable, complete, finite-volume convention, an arithmetic hyperbolic \(3\)-manifold of quaternionic origin is a quotient \(M=\Gamma\backslash\mathbb H^3\), where \(\Gamma<\operatorname{PSL}_2(\C)\) is a torsion-free lattice commensurable, up to conjugacy, with the image of \(\mathcal O^1\). Here \(k\) is a number field with exactly one complex place, \(B\) is a quaternion algebra over \(k\) ramified at every real place, \(\mathcal O\subset B\) is an order, and \(\mathcal O^1\) consists of its elements of reduced norm \(1\). The image is obtained using a splitting \(B\otimes_k\C\cong M_2(\C)\) at the complex place and projection to \(\operatorname{PSL}_2(\C)\). Two subgroups are commensurable if their intersection has finite index in each.

The supplied assertion is:
\begin{quote}
For every arithmetic hyperbolic \(3\)-manifold \(M\) of quaternionic origin, there exist a finite-sheeted covering map \(p\colon\widetilde M\to M\) and a locally trivial surface bundle \(f\colon\widetilde M\to S^1\).
\end{quote}
The arithmetic description above explicitly states a conventional interpretation; the supplied record does not itself define ``quaternionic origin'' or specify completeness, finite volume, orientability, compactness, or the permitted fiber surfaces. These conventions require confirmation against the source context. No compactness requirement on the fiber is added here.

The background describes a conditional-conjecture fragment concerning noncoherence of arithmetic lattices, whereas the supplied statement is an unconditional virtual-fibering assertion. No conditional hypothesis or noncoherence conclusion is supplied, so their intended relationship remains unresolved. The supplied label \emph{open} is a curation label, not an independently verified mathematical status.
\subsection{Short English statement}
The assertion says that every arithmetic hyperbolic three-manifold of quaternionic origin has a finite cover that is a surface bundle over a circle. The intended manifold and fiber conventions, and the relationship to the background discussion of noncoherence, require human review.
\subsection{Sources}
\begin{itemize}
\item Local UnsolvedMath registry, supplied complete record: permanent ID \(30001508\), source problem code \texttt{catalogue entry}, title \emph{Virtual Fibering of Arithmetic Hyperbolic Three-Manifolds}. No proposer is identified; the supplied proposed year is 2010.
\item Supplied source citation: \emph{Geometric Group Theory, Hyperbolic Dynamics and Symplectic Geometry} (2011), Oberwolfach Reports, pp.~1761--1825, DOI \texttt{10.4171/owr/2010/30}. \url{https://doi.org/10.4171/owr/2010/30}. The supplied record labels the statement ``exact'' and reports comparison with the official report, with notation, context, and question form normalized. This verification claim has not been independently checked here.
\item Supplied literature source, labeled ``Kapovich noncoherence theorem'': \url{https://doi.org/10.2140/gt.2013.17.39}.
\item Supplied literature source, labeled ``original OWR report'': \url{https://doi.org/10.4171/owr/2010/30}.
\item Literature checked at: \texttt{2026-08-21}. This is supplied dated curation, not independently verified. The supplied assessment labels the problem open and says that known virtual-fibration results settle several families but not the full conjectural scope represented by the extracted noncoherence fragment. The record expressly requests independent verification of time-sensitive status and cited claims.
\end{itemize}
\end{document}
